\FloatBarrier
\section{IR and Artifact Materials}
\label{sec:artifact}
The core value and index-expression forms are listed in
Tables~\ref{tab:value-expressions} and~\ref{tab:index-expressions}. Their
coordinate interpretation uses the mapping invariant in
Equation~\eqref{eq:mapping-invariant}, complete disjoint partitions from
Section~\ref{sec:mapping-algorithm}, and ordered output interfaces. Normalized
terms have the denotation in Equation~\eqref{eq:normalized-term}. Arithmetic
normalization is conditional on the selected laws and on the reduction-domain
and read-independence checks described in Section~\ref{sec:normalization}.
Algorithms~\ref{alg:mapping} and~\ref{alg:synthesis} in the main text give the construction and synthesis
procedures; neither is a claim of completeness.

The accompanying artifact should record each benchmark's definitions, guards,
explicit and synthesized constructors, arithmetic policy, and failure
classification. Reproduction materials should pin the compiler and runtime,
include independent references, and separate compilation, semantic checking,
and application measurements. \todo{Supply the versioned artifact, formal
IR grammar and complete operation denotations, exact algebraic side-condition
catalog, benchmark manifest, build and reproduction commands, environments,
and independent reference tests. These materials are not supplied in this
draft; the correctness contract and assumptions remain as stated in the main
text.}
